% Generated deterministically from the audited Markdown source.
% Responsible author: Carptopus.
\documentclass[11pt]{article}
\usepackage[a4paper,margin=27mm]{geometry}
\usepackage[T1]{fontenc}
\usepackage[utf8]{inputenc}
\usepackage{lmodern}
\usepackage{microtype}
\usepackage{amsmath,amssymb,amsthm,mathtools}
\usepackage{booktabs,tabularx,array,graphicx}
\usepackage{needspace}
\usepackage{xcolor}
\usepackage{fancyvrb}
\usepackage[colorlinks=true,linkcolor=blue!55!black,citecolor=blue!55!black,urlcolor=blue!65!black]{hyperref}
\usepackage{fancyhdr}

\newtheorem{theorem}{Theorem}[section]
\newtheorem{proposition}[theorem]{Proposition}
\newtheorem{lemma}[theorem]{Lemma}
\newtheorem{corollary}[theorem]{Corollary}
\numberwithin{equation}{section}

\pagestyle{fancy}
\fancyhf{}
\fancyhead[L]{Carptopus}
\fancyhead[R]{Tail-projection certificates}
\fancyfoot[C]{\thepage}
\setlength{\headheight}{14pt}
\setlength{\parindent}{0pt}
\setlength{\parskip}{0.45em}
\setlength{\emergencystretch}{4em}
\urlstyle{same}

\title{Tail-projection certificates for cyclotomic\\
minimum-degree polynomials beyond the reciprocal condition}
\author{Carptopus\\\href{mailto:carptopus@163.com}{carptopus@163.com}}
\date{27 August 2026}

\hypersetup{
  pdftitle={Tail-projection certificates for cyclotomic minimum-degree polynomials beyond the reciprocal condition},
  pdfauthor={Carptopus},
  pdfsubject={An explicit sufficient condition and infinite prime families for a cyclotomic nonnegative-coefficient minimum-degree problem},
  pdfkeywords={cyclotomic polynomial, nonnegative coefficients, minimum degree, separation certificate, CRT projection, infinite prime family}
}

\begin{document}
\maketitle

\begin{center}
\fcolorbox{black!35}{black!4}{\parbox{0.88\textwidth}{\centering\small
Internal candidate manuscript, version 0.1-beta. The mathematical proof and
complete Markdown manuscript have passed project-internal zero-trust audits.
External mathematical review and formal peer review are pending.}}
\end{center}

\begin{abstract}


Let \(\Phi_n(x)\) be the \(n\)-th cyclotomic polynomial, and let \(p\) be the
smallest prime divisor of \(n\). Steinberger formulated as Conjecture 1 the
assertion that the unique monic polynomial of lowest degree among the nonzero
polynomials in \(\mathbb R_{\geq0}[x]\) divisible by \(\Phi_n\) is

\nopagebreak[4]
\[
R_{n,p}(x)=1+x^{n/p}+\cdots+x^{(p-1)n/p}.
\]

He explicitly expressed skepticism about this assertion and anticipated
counterexamples. For an odd squarefree integer \(n=pq_1\cdots q_k\), where
\(p<q_1<\cdots<q_k\) and \(k\geq3\), we give an explicit sufficient
condition for this conclusion. Put \(P=q_1\cdots q_k\). For every \(d\mid P\),
let \(\rho_d\in\{1,\ldots,p-1\}\) be congruent to \(P/d\) modulo \(p\), and
write \(\sigma_d=(-1)^{k-\omega(d)}\). We prove that

\nopagebreak[4]
\[
\sum_{\substack{d\mid P,\ d<P\\ \sigma_d=-1}}d\rho_d
+p\sum_{\substack{d\mid P,\ d<P\\ \sigma_d=+1}}d
+\frac{p}{q_1-1}\sum_{\substack{d\mid P\\d<P}}d<P
\]

implies that the minimum degree is \((p-1)P\), with equality only for positive
scalar multiples of \(R_{n,p}\). The proof constructs an integer separation
certificate by applying a zero-marginal CRT tensor projection to the final
interval of \(P-1\) exponents and then cancelling its values on the regular
polygon by an explicit anchor correction.

For every fixed \(k\geq3\), the new condition holds for all sufficiently large
prime tuples

\nopagebreak[4]
\[
p<q_1<\cdots<q_k\leq(1+\delta)p,
\qquad 0<\delta<1/k.
\]

These tuples fail Steinberger's reciprocal condition
\(2/p>\sum_i1/q_i\). Consequently the theorem supplies, for each fixed number
of at least four distinct prime factors, infinitely many cases outside the
range of the previously known reciprocal sufficient condition.

\end{abstract}

\section{Introduction}


For a positive integer \(n\), consider

\nopagebreak[4]
\[
\mu(n)=
\min\{\deg f:0\ne f\in\mathbb R_{\geq0}[x],\ \Phi_n(x)\mid f(x)\}.
\tag{1.1}
\]

If \(p\) is the smallest prime divisor of \(n\), then

\nopagebreak[4]
\[
R_{n,p}(x)=1+x^{n/p}+\cdots+x^{(p-1)n/p}
\tag{1.2}
\]

is a nonnegative multiple of \(\Phi_n\), and hence

\nopagebreak[4]
\[
\mu(n)\leq (p-1)\frac np.
\tag{1.3}
\]

Steinberger \cite{steinberger2012} named equality in (1.3), together with uniqueness of the
monic minimizer, Conjecture 1, while explicitly stating that he did not
believe it in full generality and expected counterexamples. For a squarefree
factorization

\nopagebreak[4]
\[
n=pq_1\cdots q_k,
\qquad p<q_1<\cdots<q_k,
\tag{1.4}
\]

his main theorem proves the conjecture under

\nopagebreak[4]
\[
\frac2p>\sum_{i=1}^k\frac1{q_i}.
\tag{1.5}
\]

The 2012 paper left the general problem open. To the best of our knowledge,
the literature search described in Section 8 found no subsequent complete
resolution. A 2025 treatment of the corresponding integer-tiling diameter
problem still cites (1.5) as the relevant deeper sufficient theorem and
emphasizes the additional structural difficulty once three or more
nonminimal prime factors are present \cite{labazakharov2025}.

The present paper proves a different sufficient condition. Its origin is a
family of exact finite certificates developed after the first two parameter
values not covered by (1.5) were settled computationally \cite{carptopus2026}. The result
below no longer requires a linear program or a stored high-dimensional
certificate: a single interval, a tensor projection, and a closed
divisor--residue estimate provide a certificate for an infinite family.

\Needspace{0.28\textheight}
\begin{theorem}

Let \(n=pq_1\cdots q_k\) be the complete factorization of \(n\) into distinct
odd primes, with

\nopagebreak[4]
\[
p<q_1<\cdots<q_k,\qquad k\geq3,
\]

and put \(P=q_1\cdots q_k\). For every divisor \(d\mid P\), define

\nopagebreak[4]
\[
\rho_d\in\{1,\ldots,p-1\},
\qquad
\rho_d\equiv P/d\pmod p,
\tag{1.6}
\]

and

\nopagebreak[4]
\[
\sigma_d=(-1)^{k-\omega(d)},
\tag{1.7}
\]

where \(\omega(d)\) is the number of distinct prime factors of \(d\). Set

\nopagebreak[4]
\[
\mathcal B_k=
\sum_{\substack{d\mid P,\ d<P\\ \sigma_d=-1}}d\rho_d
+p\sum_{\substack{d\mid P,\ d<P\\ \sigma_d=+1}}d
+\frac{p}{q_1-1}\sum_{\substack{d\mid P\\d<P}}d.
\tag{1.8}
\]

If

\nopagebreak[4]
\[
\mathcal B_k<P,
\tag{1.9}
\]

then every nonzero \(f\in\mathbb R_{\geq0}[x]\) divisible by \(\Phi_n\)
satisfies

\nopagebreak[4]
\[
\deg f\geq(p-1)P.
\tag{1.10}
\]

Equality holds if and only if \(f=cR_{n,p}\) for some \(c>0\). In particular,
\(\mu(n)=(p-1)P\), and \(R_{n,p}\) is the unique monic minimizer.

\end{theorem}

\Needspace{0.28\textheight}
\begin{theorem}

Fix \(k\geq3\) and \(0<\delta<1/k\). Among the prime tuples satisfying

\nopagebreak[4]
\[
p<q_1<\cdots<q_k\leq(1+\delta)p,
\tag{1.11}
\]

all sufficiently large tuples satisfy (1.9), while all such tuples fail
(1.5). There are infinitely many such tuples.

In particular, if \(p_m\) denotes the \(m\)-th prime, then

\nopagebreak[4]
\[
(p,q_1,\ldots,q_k)
=(p_m,p_{m+1},\ldots,p_{m+k})
\tag{1.12}
\]

satisfies (1.9) and fails (1.5) for all sufficiently large \(m\).

\end{theorem}

Section 2 recalls the separation criterion and boundary rigidity. Sections
3--5 construct and estimate the certificate. Section 6 proves the infinite
family. Section 7 records a four-prime corollary and exact calibration, and
Section 8 explains the scope relative to earlier work.

\section{Boundary separation and rigidity}


Let \(n\) be squarefree, let \(p\) be its smallest prime divisor, and put

\nopagebreak[4]
\[
P=n/p,\qquad
H=\{0,P,\ldots,(p-1)P\},
\tag{2.1}
\]

and

\nopagebreak[4]
\[
F^+=\{0,\ldots,(p-1)P-1\}\setminus H.
\tag{2.2}
\]

Let \(V_n\subseteq\mathbb R^n\) be the coefficient-vector space of
polynomials of degree less than \(n\) divisible by \(\Phi_n\). We use the
sign-reversed form of Steinberger's separation criterion \cite[Lemma 1]{steinberger2012}.

\Needspace{0.28\textheight}
\begin{proposition}

Suppose that \(w\in V_n^\perp\) satisfies

\nopagebreak[4]
\[
w_j=0\quad(j\in H),
\qquad
w_j<0\quad(j\in F^+).
\tag{2.3}
\]

Then every nonzero \(f\in\mathbb R_{\geq0}[x]\) divisible by \(\Phi_n\) and
of degree at most \((p-1)P\) is supported on \(H\).

\end{proposition}

\begin{proof}

Pad the coefficient vector \(a\) of \(f\) with zeros to length \(n\).
Then \(a\in V_n\), so

\nopagebreak[4]
\[
0=\langle w,a\rangle
=\sum_{j=0}^{(p-1)P}w_ja_j.
\]

All summands are nonpositive. A positive coefficient at any index in \(F^+\)
would make the sum strictly negative. Hence \(a_j=0\) for every
\(j\in F^+\).
\end{proof}


\Needspace{0.28\textheight}
\begin{lemma}

If \(0\ne f\in\mathbb R_{\geq0}[x]\), \(\Phi_n\mid f\),
\(\deg f\leq(p-1)P\), and \(\mathrm{supp}(f)\subseteq H\), then

\nopagebreak[4]
\[
f(x)=cR_{n,p}(x)
\tag{2.4}
\]

for some \(c>0\).

\end{lemma}

\begin{proof}

Write \(f(x)=g(x^P)\), where \(\deg g\leq p-1\). Let \(\zeta_n\) be a
primitive \(n\)-th root. For every \(a\in(\mathbb Z/n\mathbb Z)^*\),

\nopagebreak[4]
\[
0=f(\zeta_n^a)=g((\zeta_n^P)^a).
\]

The reduction map from \((\mathbb Z/n\mathbb Z)^*\) onto
\((\mathbb Z/p\mathbb Z)^*\) is surjective by the Chinese remainder theorem.
Thus \(g\) vanishes at every primitive \(p\)-th root, so
\(\Phi_p\mid g\) in \(\mathbb R[x]\). The degree bound forces
\(g=c\Phi_p\), and nonzero nonnegative coefficients imply \(c>0\).
\end{proof}


\section{CRT projection and anchor correction}


Assume the notation of Theorem 1.1. Identify \(\mathbb Z_n\) with

\nopagebreak[4]
\[
G=\mathbb Z_p\times
\mathbb Z_{q_1}\times\cdots\times\mathbb Z_{q_k}
\tag{3.1}
\]

through the exponent bijection

\nopagebreak[4]
\[
(a,b_1,\ldots,b_k)
\longmapsto
aP+\sum_{i=1}^k b_i\frac n{q_i}\pmod n.
\tag{3.2}
\]

Under this identification, \(H\) is
\(\{(a,0,\ldots,0):a\in\mathbb Z_p\}\).
For a coordinate of length \(t\), let \(I_t\) and \(J_t\) denote the identity
and all-ones matrices. Define

\nopagebreak[4]
\[
D=\bigotimes_{t\in\{p,q_1,\ldots,q_k\}}(tI_t-J_t).
\tag{3.3}
\]

Each factor has image \(\mathbf1_t^\perp\). The standard CRT fiber
description of \(V_n\) gives

\nopagebreak[4]
\[
\mathrm{im}\,D=
\mathbf1_p^\perp\otimes
\mathbf1_{q_1}^\perp\otimes\cdots\otimes
\mathbf1_{q_k}^\perp
=V_n^\perp.
\tag{3.4}
\]

For completeness, we justify the final equality. Let \(W\) be the span of
all coordinate-fiber indicators in \(G\). A fiber in the coordinate of a
prime \(r\mid n\) is, up to a cyclic translation, the coefficient vector of

\nopagebreak[4]
\[
1+x^{n/r}+\cdots+x^{(r-1)n/r}
=\frac{x^n-1}{x^{n/r}-1},
\]

so it belongs to \(V_n\). Hence \(W\subseteq V_n\). An array is orthogonal
to every coordinate fiber exactly when every coordinate-fiber sum is zero;
therefore

\nopagebreak[4]
\[
W^\perp=
\mathbf1_p^\perp\otimes
\mathbf1_{q_1}^\perp\otimes\cdots\otimes
\mathbf1_{q_k}^\perp.
\]

The right-hand side has dimension
\((p-1)\prod_i(q_i-1)=\varphi(n)\), so
\(\dim W=n-\varphi(n)\). On the other hand, multiplication by \(\Phi_n\)
identifies the polynomials of degree less than \(n-\varphi(n)\) with \(V_n\),
and hence \(\dim V_n=n-\varphi(n)\). Thus \(W=V_n\) and the final equality
in (3.4) follows. This is also the fiber description in \cite[Section 2, pp. 5--7]{steinberger2012}.

Let

\nopagebreak[4]
\[
\mathcal I=\{n-P+1,\ldots,n-1\},
\qquad
u=D\mathbf1_{\mathcal I},
\qquad
C=\prod_{i=1}^k(q_i-1).
\tag{3.5}
\]

For \(a\in\mathbb Z_p\), put \(u_a=u(a,0,\ldots,0)\). Because the
\(p\)-coordinate factor is \(pI_p-J_p\),

\nopagebreak[4]
\[
\sum_{a\in\mathbb Z_p}u_a=0.
\tag{3.6}
\]

Define

\nopagebreak[4]
\[
Y=pC\mathbf1_{\mathcal I}
-\sum_{a\in\mathbb Z_p}
u_a\mathbf1_{\{(a,0,\ldots,0)\}},
\qquad
w=DY.
\tag{3.7}
\]

Then \(Y,w\) are integer arrays and \(w\in V_n^\perp\).

\Needspace{0.28\textheight}
\begin{lemma}

Define

\nopagebreak[4]
\[
\varepsilon(x_{q_1},\ldots,x_{q_k})
=
\prod_{i=1}^k
\frac{q_i\mathbf1_{x_{q_i}=0}-1}{q_i-1}.
\tag{3.8}
\]

Then

\nopagebreak[4]
\[
\frac{w(x)}{pC}
=u(x)-u_{x_p}\varepsilon(x_{q_1},\ldots,x_{q_k}).
\tag{3.9}
\]

Consequently \(w=0\) on \(H\). Away from \(H\),

\nopagebreak[4]
\[
|\varepsilon|\leq\frac1{q_1-1}.
\tag{3.10}
\]

\end{lemma}

\begin{proof}

For an anchor \(e_a=(a,0,\ldots,0)\),

\nopagebreak[4]
\[
D\mathbf1_{\{e_a\}}(x)
=
(p\mathbf1_{x_p=a}-1)
\prod_{i=1}^k(q_i\mathbf1_{x_{q_i}=0}-1).
\tag{3.11}
\]

Use (3.6), multiply by \(u_a\), and sum over \(a\) to obtain (3.9).
At an anchor, all factors in (3.8) equal \(1\) and \(u(x)=u_{x_p}\).
At a nonanchor, at least one final coordinate is nonzero, producing a factor
of magnitude \(1/(q_i-1)\); all other factors have magnitude at most \(1\).
This proves the two assertions.
\end{proof}


\section{Exact projection of the final interval}


For \(m\mid n\) and \(j\in\mathbb Z_n\), let

\nopagebreak[4]
\[
N_m(j)=
\#\{\ell\in\mathcal I:\ell\equiv j\pmod m\}.
\tag{4.1}
\]

Expanding (3.3) and pairing the terms indexed by \(d\) and \(pd\) gives

\nopagebreak[4]
\[
u(j)=
\sum_{d\mid P}
(-1)^{k-\omega(d)}d\,e_d(j),
\qquad
e_d(j)=pN_{pd}(j)-N_d(j).
\tag{4.2}
\]

\Needspace{0.28\textheight}
\begin{lemma}

For every \(d\mid P\),

\nopagebreak[4]
\[
N_d(j)=\frac Pd-\mathbf1_{d\mid j},
\tag{4.3}
\]

and

\nopagebreak[4]
\[
N_{pd}(j)=
\left\lfloor
\frac{P-1+(j\bmod pd)}{pd}
\right\rfloor.
\tag{4.4}
\]

If \(P/d=ap+\rho_d\), then

\nopagebreak[4]
\[
e_d(j)=
p\theta_d(j)-\rho_d+\mathbf1_{d\mid j}
\tag{4.5}
\]

for some \(\theta_d(j)\in\{0,1\}\). In particular,

\nopagebreak[4]
\[
-\rho_d\leq e_d(j)\leq p.
\tag{4.6}
\]

\end{lemma}

\begin{proof}

Write \(\ell=n-r\), where \(1\leq r\leq P-1\). Every residue modulo
\(d\) occurs \(P/d\) times in \(\{0,\ldots,P-1\}\); removing \(r=0\)
gives (4.3). Counting the positive integers \(r\leq P-1\) congruent to
\(-j\) modulo \(pd\) gives (4.4). The integer in (4.4) is either \(a\)
or \(a+1\). Substitution into \(e_d\) proves (4.5), and
\(1\leq\rho_d\leq p-1\) yields (4.6).
\end{proof}


\Needspace{0.28\textheight}
\begin{lemma}

For every \(j\in F^+\),

\nopagebreak[4]
\[
u(j)\leq
-P+
\sum_{\substack{d\mid P,\ d<P\\ \sigma_d=-1}}d\rho_d
+p\sum_{\substack{d\mid P,\ d<P\\ \sigma_d=+1}}d.
\tag{4.7}
\]

For every anchor \(h\in H\),

\nopagebreak[4]
\[
|u(h)|\leq
p\sum_{\substack{d\mid P\\d<P}}d.
\tag{4.8}
\]

\end{lemma}

\begin{proof}

For \(j\in F^+\), one has \(P\nmid j\), \(N_P(j)=1\), and
\(N_n(j)=0\). Hence \(e_P(j)=-1\), and the \(d=P\) term in (4.2)
is \(-P\). For a proper divisor with \(\sigma_d=-1\), (4.6) gives
\(\sigma_dd e_d\leq d\rho_d\); for \(\sigma_d=+1\), it gives
\(\sigma_dd e_d\leq dp\). This proves (4.7).

At an anchor the \(d=P\) term vanishes, while
\(|e_d(h)|\leq p\) for every proper divisor. Summing proves (4.8).
\end{proof}


\section{Proof of Theorem 1.1}

\begin{proof}[Proof of Theorem 1.1]

For \(j\in F^+\), Lemmas 3.1 and 4.2 give

\nopagebreak[4]
\[
\begin{aligned}
\frac{w_j}{pC}
&=u(j)-u_{j_p}\varepsilon(j_{q_1},\ldots,j_{q_k})\\
&\leq
-P+
\sum_{\substack{d\mid P,\ d<P\\ \sigma_d=-1}}d\rho_d
+p\sum_{\substack{d\mid P,\ d<P\\ \sigma_d=+1}}d\\
&\qquad
+\frac{p}{q_1-1}
\sum_{\substack{d\mid P\\d<P}}d\\
&=-P+\mathcal B_k.
\end{aligned}
\tag{5.1}
\]

Condition (1.9) makes this strictly negative. Lemma 3.1 gives \(w=0\)
on \(H\), and (3.4) gives \(w\in V_n^\perp\). Proposition 2.1 therefore
forces every nonnegative \(\Phi_n\)-multiple of degree at most
\((p-1)P\) to be supported on \(H\). Lemma 2.2 makes it a positive
scalar multiple of \(R_{n,p}\). Since \(R_{n,p}\) has degree
\((p-1)P\), the degree bound and uniqueness follow.
\end{proof}


\section{Infinite families outside the reciprocal condition}

\begin{proof}[Proof of Theorem 1.2]

Fix \(k\geq3\), \(0<\delta<1/k\), and suppose (1.11) holds. For
sufficiently large \(p\), every \(q_i<2p\). If the complement of the
prime support of \(d\mid P\) consists of one prime \(q_i\), then
\(d=P/q_i\), \(\sigma_d=-1\), and

\nopagebreak[4]
\[
\rho_d=q_i-p.
\tag{6.1}
\]

Its normalized contribution is

\nopagebreak[4]
\[
\frac{d\rho_d}{P}
=\frac{q_i-p}{q_i}
\leq\delta.
\tag{6.2}
\]

The \(k\) such terms contribute at most \(k\delta<1\). If the complement
has size \(c\geq2\), both possible upper-bound contributions are at most
\(dp\), and hence

\nopagebreak[4]
\[
\frac{\text{upper-bound contribution}}P
\leq
\frac{dp}P
=\frac{p}{P/d}
\leq\frac1{p^{c-1}}.
\tag{6.3}
\]

All complement sizes at least two therefore contribute at most

\nopagebreak[4]
\[
\sum_{c=2}^k\binom{k}{c}p^{1-c}\longrightarrow0.
\tag{6.4}
\]

The normalized anchor term is

\nopagebreak[4]
\[
\frac{p}{q_1-1}
\left(
\prod_{i=1}^k\left(1+\frac1{q_i}\right)-1
\right)
\longrightarrow0.
\tag{6.5}
\]

Combining the three estimates gives the uniform bound, over all prime tuples
satisfying (1.11),

\nopagebreak[4]
\[
\frac{\mathcal B_k}{P}
\leq
k\delta+
\sum_{c=2}^k\binom{k}{c}p^{1-c}
+\frac{p}{q_1-1}
\left(\left(1+\frac1p\right)^k-1\right)
=k\delta+O_k(p^{-1}).
\tag{6.6}
\]

Since \(k\delta<1\), condition (1.9) eventually holds uniformly. Conversely,

\nopagebreak[4]
\[
p\sum_{i=1}^k\frac1{q_i}
\geq\frac{k}{1+\delta}>2,
\tag{6.7}
\]

so Steinberger's condition (1.5) fails.

The prime number theorem implies that
\((p,(1+\delta)p]\) contains at least \(k\) primes for every sufficiently
large prime \(p\), producing infinitely many tuples. Also
\(p_{m+i}/p_m\to1\) for fixed \(i\), which proves the consecutive-prime
statement.
\end{proof}


\section{Four-prime corollary and exact calibration}


For \(k=3\), write \((q_1,q_2,q_3)=(q,r,s)\).

\Needspace{0.28\textheight}
\begin{corollary}

Let \(n=pqrs\) be a product of distinct odd primes
\(p<q<r<s<2p\). If

\nopagebreak[4]
\[
\begin{aligned}
&qr(s-p)+qs(r-p)+rs(q-p)+p(q+r+s)+p\\
&\qquad
+\frac{p}{q-1}(qr+qs+rs+q+r+s+1)
<qrs,
\end{aligned}
\tag{7.1}
\]

then the conclusion of Theorem 1.1 holds.

\end{corollary}

\begin{proof}

For \(d=qr,qs,rs\), the residues \(\rho_d\) are respectively
\(s-p,r-p,q-p\). For \(d=1\), use \(\rho_1\leq p\). The other
proper-divisor terms agree with the bounds in (7.1), so its left side
majorizes \(\mathcal B_3\).
\end{proof}


The exact arithmetic checker gives the following consecutive-prime examples:

\begin{center}
\small
\resizebox{\textwidth}{!}{%
\begin{tabular}{l l l l}
\toprule
Total prime factors & Prime tuple & \(n\) & \(\mathcal B_k/P\) \\
\midrule
4 & \(23,29,31,37\) & \(765049\) & \(31795/32116\) \\
5 & \(37,41,43,47,53\) & \(162490421\) & \(165666531/175665320\) \\
6 & \(59,61,67,71,73,79\) & \(98733594781\) & \(96500955779/100407045540\) \\
7 & \(97,101,103,107,109,113,127\) & \(168897325606883\) & \(149066984374577/174120954233900\) \\
\bottomrule
\end{tabular}%
}
\end{center}


Every row satisfies \(\mathcal B_k/P<1\) and fails (1.5). These examples
calibrate the formulas and are not used in the infinite-family proof.

\section{Scope and relation to earlier work}


Steinberger \cite{steinberger2012} introduced the problem, the separation criterion, and the
reciprocal condition (1.5). Theorem 1.1 uses the same separation framework
but a different explicit certificate; Theorem 1.2 reaches infinitely many
tuples outside (1.5).

The companion fixed-parameter paper \cite{carptopus2026} proves the conjectured minimum and
uniqueness for

\nopagebreak[4]
\[
46189=11\cdot13\cdot17\cdot19
\quad\text{and}\quad
96577=13\cdot17\cdot19\cdot23
\]

using exact stored integer certificates. Those small tuples need not satisfy
(1.9), so the present theorem does not logically subsume the fixed-parameter
results. The fixed certificates motivated the projection construction,
whereas the present work extracts a symbolic mechanism for infinitely many
additional cases.

Reddy \cite{reddy2011} treats lowest-degree \(0,1\)-polynomials and minimal vanishing sums
for even orders with at most three distinct prime factors. Kepka and
Korbel\'a\v{r} \cite{kepkakorbelar2012} treat bounds for nonnegative multiples of a general polynomial
and determine the cubic case. Neither directly supplies (1.8), the
anchor-corrected CRT projection, or Theorem 1.2.

This paper does not prove or disprove Steinberger's conjecture in full. It
does not claim that (1.9) is necessary or that the final-interval ansatz works
for every prime tuple. Failure of (1.9) is only failure of this sufficient
estimate, not evidence of a counterexample.

The novelty and open-status check was refreshed on 27 August 2026. We
searched arXiv, OpenAlex, Semantic Scholar, and general scholarly web indexes
using the title and DOI of \cite{steinberger2012}, the exact minimum-degree problem, and
combinations of ``cyclotomic'', ``nonnegative coefficients'', ``tail projection'',
``divisor residue'', and ``separation certificate''. We also inspected the
forward-citation neighborhood of \cite{steinberger2012} and the 2025 journal version of \cite{labazakharov2025}.
We found no source directly stating (1.8), the anchor-corrected final-interval
projection, or Theorem 1.2. This is a bounded literature search, not a proof
of global priority; differently phrased, unindexed, or concurrent work may
still narrow the contribution.

\section{Reproducibility}


The proof is symbolic and does not depend on finite enumeration or
floating-point optimization. The accompanying checker
{\small\path{verification/verify_general_tail_condition.py}} uses exact
integer and rational arithmetic to verify primality and factor ordering,
evaluate every proper-divisor term in \(\mathcal B_k\), test (1.9), and test
Steinberger's reciprocal condition. For example:

\begin{Verbatim}[fontsize=\scriptsize]
python -X utf8 `
  'verification\verify_general_tail_condition.py' `
  --primes 37 41 43 47 53
\end{Verbatim}


The script rejects inputs outside the theorem's domain, including even,
composite, repeated, or unordered factors. It is a calibration aid, not a
substitute for the proof.

\section{AI-assisted research and writing disclosure}


Computational exploration, proof organization, adversarial checking, and
language editing were assisted by OpenAI Codex. Carptopus selected and
reviewed the mathematical claims, verified the retained evidence, and
assumes responsibility for the manuscript.

\begin{thebibliography}{9}
\footnotesize
\raggedright

\bibitem{steinberger2012}
J.~P. Steinberger,
\emph{The lowest-degree polynomial with nonnegative coefficients divisible by
the $n$-th cyclotomic polynomial},
Electronic Journal of Combinatorics \textbf{19} (2012), no.~4, Paper P1.
\href{https://doi.org/10.37236/2755}{doi:10.37236/2755}.

\bibitem{labazakharov2025}
I.~\L aba and D.~Zakharov,
\emph{On the minimal period of integer tilings},
Bulletin of the London Mathematical Society \textbf{57} (2025), no.~4,
1160--1170.
\href{https://doi.org/10.1112/blms.70023}{doi:10.1112/blms.70023}.

\bibitem{carptopus2026}
Carptopus,
\emph{Exact certificates for the first two cases not covered by Steinberger's
cyclotomic minimum-degree theorem},
Public Beta v0.1-beta (2026).
\href{https://doi.org/10.5281/zenodo.22124096}{doi:10.5281/zenodo.22124096}.

\bibitem{reddy2011}
A.~S. Reddy,
\emph{The lowest degree $0,1$-polynomial divisible by cyclotomic polynomial},
arXiv:1106.1271v2 (2011).
\href{https://arxiv.org/abs/1106.1271}{arXiv:1106.1271}.

\bibitem{kepkakorbelar2012}
T.~Kepka and M.~Korbel\'a\v{r},
\emph{The lowest-degree polynomials with non-negative coefficients},
arXiv:1210.6868 (2012).
\href{https://arxiv.org/abs/1210.6868}{arXiv:1210.6868}.

\end{thebibliography}

\end{document}
